\documentclass[11pt]{article}

\usepackage[margin=1in]{geometry}
\usepackage[T1]{fontenc}
\usepackage[utf8]{inputenc}
\usepackage{lmodern}
\usepackage{microtype}
\usepackage{amsmath, amssymb, amsthm, mathtools}
\usepackage{bm}
\usepackage{enumitem}
\usepackage{graphicx}
\usepackage{hyperref}
\usepackage{xcolor}
\usepackage{xypic}
\usepackage[all,line,color]{xy}
\usepackage{graphicx}
\usepackage{pagecolor}
\pagecolor{white}
\usepackage{hyperref}
\usepackage{wrapfig
}\usepackage{caption}
\usepackage{tikz}
\usetikzlibrary{arrows.meta}
\usepackage{amsmath,amssymb}
\usepackage{algorithm}
\usepackage{algpseudocode}

% Define the custom inline macro
\DeclareRobustCommand{\externallink}{%
  \tikz[
    x=1ex,
    y=1ex,
    baseline=-0.05ex,
    line width=0.18ex,
    line cap=round,
    line join=round
  ]{%
    % outer box, intentionally open at top-right
    \draw (0,0) -- (0,1.05) -- (0.68,1.05);
    \draw (0,0) -- (1.05,0) -- (1.05,0.68);

    % arrow corner
    \draw (0.78,1.25) -- (1.25,1.25) -- (1.25,0.78);

    % arrow shaft
    \draw[-{Stealth[length=0.32ex,width=0.28ex]}]
      (0.48,0.48) -- (1.28,1.28);
  }%
}

\captionsetup[figure]{
  format=plain,
  justification=justified,
  singlelinecheck=false,
  margin=1.75cm,
  font=small
}

\hypersetup{
    colorlinks=true,
    linkcolor=blue!50!black,
    citecolor=blue!50!black,
    urlcolor=blue!50!black,
    pdftitle={Problem agnostic hierarchies, path-volume, and log speed-up in RL},
    pdfauthor={Tyler Foster}
}

\author{Tyler Foster}
\date{\today}

\newtheorem{definition}[subsubsection]{Definition}
\newtheorem{assumption}[subsubsection]{Assumption}
\newtheorem{assumptions}[subsubsection]{Assumptions}
\newtheorem{construction}[subsubsection]{Construction}
\newtheorem{example}[subsubsection]{Example}
\newtheorem{remark}[subsubsection]{Remark}
\newtheorem{ermark}[subsubsection]{}
\newtheorem{warning}[subsubsection]{Warning}
\newtheorem{question}[]{Question}
\newtheorem{theorem}[subsubsection]{Theorem}
\newtheorem{proposition}[subsubsection]{Proposition}
\newtheorem{lemma}[subsubsection]{Lemma}
\newtheorem{corollary}[subsubsection]{Corollary}
\newtheorem{principle}[subsubsection]{Principle}
\newtheorem{observation}[subsubsection]{Observation}
\newtheorem{run_analysis}[subsubsection]{Runtime analysis}


\newcommand{\Act}{\mathsf{A}}
\newcommand{\src}{\operatorname{src}}
\newcommand{\tgt}{\operatorname{tgt}}
\newcommand{\Path}{\operatorname{Path}}
\newcommand{\Hist}{\operatorname{Hist}}
\newcommand{\supp}{\operatorname{supp}}
\newcommand{\Vol}{\operatorname{Vol}}
\newcommand{\E}{\mathbb{E}}
\newcommand{\R}{\mathbb{R}}
\newcommand{\N}{\mathbb{N}}
\newcommand{\PVol}{\mathcal{P}\text{-}\mathrm{Vol}}
\newcommand{\TV}{\operatorname{TVol}}
\newcommand{\diam}{\operatorname{diam}}
\newcommand{\im}{\operatorname{im}}
\newcommand{\id}{\operatorname{id}}
\newcommand{\cost}{\operatorname{Cost}}
\newcommand{\fiber}{\operatorname{Fib}}
\newcommand{\quot}{\mathbin{/}}
\newcommand{\localfigdir}{assets/images/mathematical_model}
\newcommand{\projectfigdir}{../../assets/images/mathematical_model}

% Prefer local bundled figures when this draft is compiled from the zip.
% If this file is dropped into docs/design/ inside the repository, fall back
% to the repository's existing assets directory. If neither exists, preserve
% the figure slot as an explicit placeholder rather than silently deleting it.
\newcommand{\safeincludegraphics}[2][]{%
  \IfFileExists{\localfigdir/#2}{%
    \includegraphics[#1]{\localfigdir/#2}%
  }{%
    \IfFileExists{\projectfigdir/#2}{%
      \includegraphics[#1]{\projectfigdir/#2}%
    }{%
      \fbox{\begin{minipage}[c][1.55in][c]{0.82\linewidth}
        \centering\small Figure file preserved but not found in this build:\\
        \texttt{#2}
      \end{minipage}}%
    }%
  }%
}

\begin{document}


Ok I see a tiny part of what's happening. It's only a very small piece of the puzzle. It's maybe the simplest non-trivial case, and it only says something about the underlying topology/geometry. It doesn't clarify anything about the coordinate rings or structure sheaves living on these spaces, which is to say that iiit doesn't say much about the log semirings living over / guiding the topology in the general case, but maybe it start to hint at where/what they might be.

\vskip .2cm

\begin{subsection}{Main observation}

So here's what I noticed. I'm just stating it succinctly. You'll have to think through details:

\begin{observation}
\normalfont
Consider the processes of setting up and running binary search on an  ordered set $\mathcal{S}$ containing $2^{N}$ elements, i.e., on a \texttt{list} of length $2^{N}$.

We know that there are two related but very different structures we can build on $\mathcal{S}$ to run binary search:
\begin{enumerate}
\item\vskip -.2cm
{\em \cite{MalikThesis}-style quotient tower}.
\item\vskip -.2cm
{\em Binary search tree}
\vskip -.2cm
\end{enumerate}
The apparent duality between these two objects is actually much deeper than just two possible choices, and is more like \href{}{{\em Gelfand duality}}, that is, duality between spaces and rings via $\text{Spec}$:
\begin{enumerate}
\item
{\em The \cite{MalikThesis}-style quotient tower} is a finite piece of the inverse system of discrete sets that induces the profinite, inverse limit topology on the ring of $2$-adic (or {\em dyadic}) integers $\widehat{\mathbb{\ Z\!}_{\ \!2}}$. For this particular problem of binary search, our quotient tower identifies our set o nodes $\mathcal{S}$ itself with the quotient $\mathbb{Z}/2^{N}\mathbb{Z}$:
$$
\begin{xy}
(0,0)*+{
\widehat{\mathbb{\ Z\!}_{\ \!2}}
}="1";
(20,0)*+{
\mathbb{Z}/2^{N}\mathbb{Z}
}="2";
(32,0)*+{
\cong\;\mathcal{S}
}="3";
{\ar@{->>} "1"; "2"}
\end{xy}
$$
\item
{\em The binary search tree} is a piece  of the Berkovich analytic space associated to $\widehat{\mathbb{\ Z\!}_{\ \!2}}$, the space denoted ``$\mathcal{M}(\widehat{\mathbb{\ Z\!}_{\ \!2}})$ in the Berkovich literature. The Berkovich space {\em is} the space being topologized by the inverse system quotient tower.
\end{enumerate}
\end{observation}

So the picture that is emerging is one where, for binary search at least, we're equipping a list we need to search through with the topology of a Berkovich space so that we can move through in logarithmic time.

So the {em vague} general picture I have is one where we're doing something {\em like} a synthetic ``Berkovich-like topologization'' of our graph so that the graph's nodes embed into a dual, hidden ``Berkovich-like'' space (not sure how much such a thing actually exists in our case) that we traverse really quickly.\footnote{\ Just an aside: {\em duality} is one of the most important things going on in tropical geometry: \url{https://pubs.ams.org/journals/jams/2018-31-02/S0894-0347-2017-00890-7}}

\hfill\textemdash\ Tyler


\vskip .2cm



\end{subsection}

\begin{subsection}{A bit of a mess below, but lots of important pieces}

The guiding algebro-geometric example here is $\mathbb{G}_{\mathrm{m}}$

Below is some output from GPT 5.5 that clarifies the role of the logarithm in all this. The drawings below are mine, except the Lionel Lang image, which comes from a paper that's not mine.

\end{subsection}

\vskip .5cm

\hrule

\vskip .5cm

The shared concrete mechanism is:

\[
\text{multiplicative scale}
\quad\longrightarrow\quad
\text{additive depth},
\]
\[
\text{fine metric detail}
\quad\longrightarrow\quad
\text{coarse order or valuation data},
\]
\[
\text{ordinary geometry}
\quad\longrightarrow\quad
\text{tree, fan, quotient skeleton}.
\]

For tropicalization with \(0<b<1\), the meaningful expression is not a fixed
\(\log_b(x)\) as \(b\to 0\), since fixed \(x\) mostly degenerates. The meaningful
move is to write

\[
x(b) = b^a,
\]

so that

\[
\log_b x(b) = a.
\]

\begin{figure}[!ht]
    \vskip -.5cm
    $$
    \!\!\!\!\!\!\!\!\!\!\!\!\!\!\!\!\!\!\!\!
    \includegraphics[scale=0.28]{assets/images/Ctimes.png}
    $$
    \vskip -.5cm
    \caption{[...]}
    \label{fig:Ctimes}
\end{figure}

For sums, one obtains

\[
\log_b\!\left(b^a + b^c\right) \longrightarrow \min(a,c)
\qquad \text{as } b\to 0.
\]

\begin{figure}[!ht]
    \vskip -.5cm
    $$
    \!\!\!\!\!\!\!\!\!\!\!\!\!\!\!\!\!\!\!\!
    \includegraphics[scale=0.78]{assets/images/lionel_lang.png}
    $$
    \vskip -.5cm
    \caption{From a paper by Lionel Lang: Applying $\text{log}_{\ }|-|$ to curves in a toric variety tend to produce tropical-looking things that aren't quite polyhedral, and look sort of like biological amoebas posing as polyhedral geometry.}
    \label{fig:lionel_lang}
\end{figure}

Thus tropicalization says:

\[
\text{keep the exponent or scale;}
\]
\[
\text{throw away coefficient-level detail;}
\]
\[
\text{turn addition into } \min;
\]
\[
\text{turn multiplication into addition.}
\]


\begin{figure}[!ht]
    \vskip -.5cm
    $$
    \includegraphics[scale=0.28]{assets/images/Qptimes.png}
    $$
    \vskip -.5cm
    \caption{[...]}
    \label{fig:Qptimes}
\end{figure}


Now compare the binary contraction tower. Take a dyadic grid. At depth \(k\),
cells have size

\[
2^{-k}.
\]

A binary quotient or coset tier forgets all detail below that scale. A point is
no longer ``this exact grid point''; it is ``this dyadic coset at depth \(k\).''
The important invariant becomes

\[
\text{scale} = 2^{-k},
\]

and therefore

\[
\text{depth} = -\log_2(\text{scale}) = k.
\]

So the binary tower is also doing:

\[
\text{keep the exponent or depth;}
\]
\[
\text{throw away within-cell detail.}
\]

That is the same move as tropicalization.

The deep reason these feel like the same phenomenon is that both are valuation
machines. A valuation converts size or order into an additive number:

\[
v(xy) = v(x) + v(y),
\]

and

\[
v(x+y) \geq \min(v(x),v(y)).
\]

In the tropical limit, this inequality becomes the governing rule:

\[
v(x+y) = \min(v(x),v(y))
\]

generically.

In the binary coset tower, the valuation is dyadic depth:

\[
v_2(\text{scale}) = \text{depth}.
\]

Equivalently, for two grid points,

\[
\text{distance scale} \sim 2^{-k}
\qquad\Longleftrightarrow\qquad
\text{separation depth} \sim k.
\]

Two points are indistinguishable until the tier where their binary addresses
split. The tower keeps that split-depth and discards the lower bits. That is a
tropical or ultrametric kind of geometry.

So a useful explanation is:

\begin{quote}
Tropicalization and binary contraction towers are both ways of replacing metric
magnitude by scale depth. The logarithm is not an extra analytic trick; it is
the coordinate change from multiplicative refinement to additive tier index.
\end{quote}

Or, more compactly:

\[
\text{tropicalization:}\qquad x=b^a \longmapsto a,
\]

\[
\text{binary tower:}\qquad \text{cell size}=2^{-k} \longmapsto k.
\]

Both say: the real object is not the tiny quantity itself, but the exponent at
which it lives.

That is why \(\lim_{b\to 0}\log_b\) and a tower of binary cosets can be
understood as manifestations of the same idea: each builds a limiting skeleton
where only dominance, scale, and ancestry survive.





\bibliographystyle{alpha}
\bibliography{logHRL}


\end{document}
